\begin{afterword}
\summaryhead

The post returns to Keats's ``dull catalogue of common things'' and asks what goes into it. The answer is whatever is known, knowable, lawful: in a word, real. The rainbow was demoted for the sin of having a scientific explanation, and the author suspects that Keats did not know the explanation himself, and that merely being told someone else knew was too much.

Since mystery is a fact about the mind, not about the world, everything that exists is liable to end up in the catalogue sooner or later. So the choice is either to decide that real, explicable things are worth caring about, or to suffer an ennui that cannot be resolved; self-deception is ruled out. Seen this way, scientists, who become fascinated by pocket lint and rainbows, are the people able, in principle, to enjoy life in the real universe.

\responsehead

This short post gives little ground for criticism. Its argument is brief and follows from its premises. If a thing becomes dull once it is known or knowable, and nothing real is mysterious in itself, then that standard will make everything real dull in time, and whoever holds it will be disappointed. The post says ``liable'' and ``sooner or later,'' and does not overstate the conclusion.

Its reading of Keats takes the ``dull catalogue'' at its word. The passage says the rainbow was catalogued because ``We know her woof, her texture,'' and it complains that philosophy will ``Conquer all mysteries by rule and line.'' This is the complaint the lines make, and here the post answers it directly; the notes on ``Explaining vs.\ Explaining Away'' observe that the earlier post had answered a different one.

What the post says about Keats himself is offered as suspicion, three times. The only evidence given in these posts for what Keats knew is the toast discussed in the notes on ``Fake Reductionism.''

In short: a brief and sound argument that valuing only the unexplained leads to disappointment, with suspicions about Keats that are marked as suspicions.
\end{afterword}
